% Generated from paper/source. Do not edit by hand.
\section{Interpretation of learned state}
\subsection{What exact equivalence permits us to say}
Complex coordinates can provide an economical language for rotations, relative phase, and ordered evidence updates. They may make a mechanism easier to write or constrain. In Q-Neuro, trajectories, amplitudes, entropies, state velocities, and phase-sensitive readouts are observable, and destructive interventions show dependence on several of these quantities. Exact-real equivalence means that such observables should be interpreted as coordinates of a structured dynamical system, not as evidence that the system requires intrinsically complex arithmetic.

\subsection{Probe evidence is not unique}
Earlier frozen-state probes recovered synthetic simulator factors such as mechanism, localization, temporality, and context from complex states. Conventional recurrent and state-space controls recovered the same factors as well or better. Linear decodability establishes accessibility to a probe, not causal use, disentanglement, or human interpretability. The current paper therefore treats probe results as historical mechanism observations rather than support for the primary claim.

\subsection{A more defensible explanatory object}
The strongest explanation produced by the project is the comparator transformation itself. The positive result versus the two-channel model is real within its declared benchmark. The exact block shows which part is a coordinate choice. The best-real envelope shows that the claimed advantage is not stable under control strengthening. The law failure shows that a compact post hoc pattern does not transfer. These linked artifacts explain why the scientific conclusion changed without erasing the original evidence.

\subsection{Limits of mechanistic language}
Terms such as coherence, interference, Hamiltonian, and measurement are mathematical descriptions in this repository. They do not establish a physical quantum substrate, biological plausibility, or a cognitive theory. A future mechanistic claim would require interventions that discriminate between mapped parameterizations, predeclared causal predictions, and independent replication on tasks whose generators were not designed by the project.
